\chapter{Axios and the Frontend API Layer}

\section{A Centralized Axios Instance}

Every service file imports one shared, pre-configured axios instance rather
than calling \texttt{axios.get(...)} with a hardcoded URL.

\begin{lstlisting}[language=JavaScript]
// src/api/axios.js
import axios from "axios";

const api = axios.create({
  baseURL: import.meta.env.VITE_API_URL, // e.g. http://localhost:5000/api
  withCredentials: true,                 // send the JWT cookie automatically
  headers: { "Content-Type": "application/json" },
});

export default api;
\end{lstlisting}

\begin{itemize}
  \item \textbf{baseURL} -- \texttt{api.get("/tasks")} hits \texttt{\{VITE\_API\_URL\}/tasks} automatically.
  \item \textbf{headers} -- defaults applied to every request, overridable per call.
  \item \textbf{withCredentials} -- required for cookie-based JWT auth (Chapter 14).
\end{itemize}

\begin{notebox}[NOTE]
\texttt{import.meta.env.VITE\_API\_URL} comes from Vite's env system; define it
in \texttt{frontend/.env} -- Vite only exposes \texttt{VITE\_}-prefixed variables.
\end{notebox}

\section{Interceptors}

Interceptors run code on every request or response without repeating it in
each service function.

\begin{lstlisting}[language=JavaScript]
// src/api/axios.js (continued)
import { store } from "../store"; // or however you track loading globally

api.interceptors.request.use((config) => {
  store.setLoading(true);
  return config;
});

api.interceptors.response.use(
  (response) => {
    store.setLoading(false);
    return response;
  },
  (error) => {
    store.setLoading(false);

    if (error.response?.status === 401) {
      window.location.href = "/login"; // session expired or invalid
    }

    return Promise.reject(error);
  }
);
\end{lstlisting}

This gives every screen a global loading indicator and one place to handle
"logged out" -- components don't need to check for 401 themselves.

\section{Error Handling in Service Functions}

\begin{lstlisting}[language=JavaScript]
// services/taskService.js
import api from "../api/axios";

export async function getTasks() {
  try {
    const res = await api.get("/tasks");
    return res.data.data; // unwrap { success, data }
  } catch (err) {
    const message = err.response?.data?.message || "Failed to load tasks";
    throw new Error(message);
  }
}
\end{lstlisting}

\begin{lstlisting}[language=JavaScript]
// pages/Dashboard.jsx (usage)
useEffect(() => {
  getTasks()
    .then(setTasks)
    .catch((err) => setError(err.message));
}, []);
\end{lstlisting}

\begin{tipbox}[TIP]
Keep \texttt{try/catch} inside the service function -- components should only
handle the friendly error message it throws, not raw axios errors.
\end{tipbox}

\begin{warnbox}[WARNING]
Don't swallow errors silently (\texttt{catch(() => \{\})}); surface or log them,
or failures become invisible.
\end{warnbox}
